\section{Introduction}\label{sec:intro}

A thousand-rupee note and a sheet of paper printed with the same inks weigh the same. The note buys
groceries and the sheet buys nothing, and the difference lies in the shopkeeper, the bank that issued
the note and everyone who will take it tomorrow because they expect others to take it too. Social
ontology studies differences of this kind: how money, marriages, licenses, committees and decisions
come to exist because people take them to exist \citep{searle1995}. Much of that taking is now done
through records, and the records are moving out of drawers and registries into systems that hold
them for us: platforms, clouds, shared ledgers, and software agents that write records on our
behalf. Each of these systems makes a choice, usually without saying so, about who settles which
facts. A platform that holds everyone's records settles all of them. A global ledger makes every
participant agree on every fact.

This paper asks which social facts need an authority behind them, someone with the standing to
settle them, and which ones anybody holding the right records can settle alone. The answer we defend
turns on whether a fact asks for something to exist or for something not to exist.

The question comes out of earlier work on collaborations among parties that keep their own records.
\emph{Collaboration Is Composition} \citep{karna2026} models each party as a machine over an
append-only log of signed records, and a collaboration as such machines wired together, so that a
collaboration is itself a machine and can be wired into a larger one. It proves, in a form of the
CALM theorem \citep{hellerstein2020,ameloot2013}, that when every machine only adds to what it has,
the parties reach the same logs and views whatever order their messages arrive in, with no
orchestrator keeping them in step. It also proves the converse case: a step that depends on a record
being absent can come out differently under two fair orders of delivery, and there the model asks
for a co-signature, or a timestamp from a witness the parties share. That paper says where
coordination is needed in a collaboration of machines. It leaves open what the coordinated steps are
in the social world the records describe, and who may take them. This paper takes up both questions.

Our thesis is that the CALM line runs through social facts as soon as those facts are grounded in
records kept in many hands. A fact whose grounding condition only asks that certain records exist
can be settled by anyone who holds them, in any order, and never has to be taken back. A fact whose
grounding condition asks for an absence cannot be settled that way by anyone, however careful, until
some coordination has happened. We call the coordinating act a seal: a declaration, in Searle's
sense, that a set of records counts as complete for a question. Once a seal exists the absence fact
rests on the seal, which is a record like any other, and the fact travels freely again. A seal can
come from a chair, a court or a registry, and it can also come from the parties the question
concerns, signing together in the way Gilbert describes joint commitment. Authority is needed where
absence matters, and where it is needed it can be shared.

A seal settles a question for everyone who accepts it, and whoever holds the standing to seal can
seal in bad faith. The second half of the paper therefore asks who may seal. We argue that the rules
which confer that standing have to be written as records too, as Hart's secondary rules, built by the
people they bind through consensus, customized by groups through additions that stay monotone, and
changed in any other way only through the legal pathways for changing regulation. Governing bodies
act through one facade over a quorum of their members, and appeals end in institutions outside the
collaboration.

The paper proceeds as follows. \secref{sec:collab} sets out the setting of monotone, composable
collaborations. \secref{sec:grounding} introduces grounding conditions, and \secref{sec:sort} sorts
them. \secref{sec:seals} defines seals and shows how they can be shared. \secref{sec:anchoring}
asks who may seal. \secref{sec:objections} answers the objections a philosopher would raise, and \secref{sec:conclusion}
closes. Throughout, the claim is conditional: insofar as an institution's facts are functions of
records \citep{ferraris2013}, the sort applies to them.
